\documentclass[border=6pt]{standalone}
\usepackage[UTF8,fontset=windows]{ctex}
\usepackage{chemfig}
\usepackage{tikz}
\usetikzlibrary{arrows.meta,calc,positioning}

\tikzset{
  molecule/.style={
    align=center,
    text=black,
    inner sep=2pt,
    font=\sffamily\footnotesize
  },
  reaction/.style={
    -{Latex[length=2.0mm,width=1.35mm]},
    draw=black,
    line width=0.5pt
  },
  route label/.style={
    fill=white,
    text=black,
    inner sep=1.2pt,
    font=\sffamily\scriptsize
  }
}

\begin{document}
\sffamily
\setchemfig{atom sep=2.0em,bond style={line width=0.5pt,draw=black}}
\begin{tikzpicture}[node distance=12mm and 22mm]
  \node[molecule] (ethanol) {
    \chemfig{CH_3-CH_2-OH}\\[-1pt]
    乙醇
  };

  \node[molecule, right=27mm of ethanol] (acetaldehyde) {
    \chemfig{CH_3-C(=[2]O)-H}\\[-1pt]
    乙醛
  };

  \node[molecule, right=30mm of acetaldehyde] (butadiene) {
    \chemfig{CH_2=CH-CH=CH_2}\\[-1pt]
    丁二烯（C4烯烃示例）
  };

  \node[molecule, below=18mm of acetaldehyde] (butanol) {
    \chemfig{CH_3-CH_2-CH_2-CH_2-OH}\\[-1pt]
    正丁醇（高碳醇示例）
  };

  \node[molecule, below=18mm of ethanol] (ethylene) {
    \chemfig{CH_2=CH_2}\\[-1pt]
    乙烯
  };

  \draw[reaction] (ethanol) -- node[route label,above]{脱氢} (acetaldehyde);
  \draw[reaction] (acetaldehyde) -- node[route label,above]{多步转化} (butadiene);
  \draw[reaction] (ethanol) -- node[route label,left]{脱水} (ethylene);
  \draw[reaction] (ethanol.south east) -- node[route label,sloped,above]{碳链增长} (butanol.north west);
\end{tikzpicture}
\end{document}
